\documentclass[11pt]{article}

\usepackage[T1]{fontenc}
\usepackage[utf8]{inputenc}
\usepackage{amsmath,amssymb,amsthm,mathtools}
\usepackage{microtype}
\usepackage[a4paper,margin=30mm]{geometry}
\usepackage{enumitem}
\usepackage{booktabs}
\usepackage{hyperref}
\usepackage[nameinlink,noabbrev]{cleveref}

\hypersetup{
  colorlinks=true,
  linkcolor=blue!55!black,
  citecolor=blue!55!black,
  urlcolor=blue!55!black,
  pdftitle={The source horizon of a meromorphic divisor},
  pdfauthor={Michael Haufschild}
}

\newtheorem{theorem}{Theorem}[section]
\newtheorem{proposition}[theorem]{Proposition}
\newtheorem{corollary}[theorem]{Corollary}
\newtheorem{lemma}[theorem]{Lemma}
\theoremstyle{definition}
\newtheorem{definition}[theorem]{Definition}
\theoremstyle{remark}
\newtheorem{remark}[theorem]{Remark}

\newcommand{\C}{\mathbb C}
\newcommand{\R}{\mathbb R}
\newcommand{\N}{\mathbb N}
\newcommand{\Res}{\operatorname*{Res}}
\newcommand{\Pois}{\operatorname{Pois}}
\newcommand{\ord}{\operatorname{ord}}
\newcommand{\Rea}{\operatorname{Re}}
\newcommand{\Ima}{\operatorname{Im}}
\newcommand{\cR}{\mathcal R}
\newcommand{\cE}{\mathcal E}
\newcommand{\cW}{\mathcal W}
\newcommand{\dd}{\,\mathrm d}

\title{\textbf{The source horizon of a meromorphic divisor}\\
\large An exact information law for escaping Riemann zeros}
\author{Michael Haufschild}
\date{}

\begin{document}

\maketitle

\begin{abstract}
A divisor can escape every fixed compact observer without losing its information.  We prove the
exact law governing where that information goes.

Let \(F\) be meromorphic, let \(c\in\C\) be a marked source, and let \(\rho\ne c\) be a divisor
point of order \(m\).  Multiplication of the logarithmic response \(F'/F\) by
\(\exp(x/(s-c))\) before residue extraction transports \(m\) into the coherent mode
\[
  M_{\rho,c}(x)=m\exp\!\left(\frac{x}{\rho-c}\right).
\]
At the canonical observation time \(x=|\rho-c|^2\tau\), reciprocal displacement becomes direct
conjugate displacement,
\[
  M_{\rho,c}(|\rho-c|^2\tau)
    =m e^{\overline{\rho-c}\tau},
\]
and every factorially normalized derivative channel obeys
\[
 \frac{x^r}{r!}\bigl|\partial_x^rM_{\rho,c}(x)\bigr|^2
 =|m|^2e^{[2\Rea(\rho-c)+1]\tau}
   \Pr\{\Pois(\tau)=r\}.
\]
Thus each source has a neutral horizon
\(\Rea\rho=\Rea c-\tfrac12\).  Radial distance has become observation time, imaginary
displacement has become phase frequency, and signed distance from the horizon is the only
surviving exponential rate.

For the completed Riemann response the arithmetic source is \(c=1\), so the neutral horizon is
exactly the critical line.  The same mode is the Poisson completion of the inverse Li--Cayley
orbit.  Its packet drift and growth rate reconstruct the zero up to conjugation, and every
strict-right zero is a positive Lyapunov mode.  We then couple this law to an explicit sequence of
functional-equation-preserving, hypertranscendental quartet perturbations.  Every holomorphic jet
converges compact-open to \(\xi\), the complete rational differential matroid is unchanged, and
the nonnegative Ward density vanishes on every bounded window while retaining mass \(8\pi\) per
quartet.  Canonical source observation nevertheless restores an exact height-independent packet.

The result does not prove the Riemann hypothesis.  It gives a closed explanation of a recurring
obstruction in analytic, Hilbert-space, differential-algebraic, and index approaches: fixed-scale
observables lose defects at the noncompact boundary, whereas the source-selected inverse-drift
scale conserves their full oriented divisor data.
\end{abstract}

\medskip
\noindent\textbf{Keywords.}
Riemann hypothesis; meromorphic divisor; logarithmic derivative; Li coefficients; coherent state;
Poisson law; hypertranscendence; Ward index; compact-open convergence.

\medskip
\noindent\textbf{2020 Mathematics Subject Classification.}
Primary 11M26, 30D30; Secondary 11J91, 46E20, 12H05.

\section{Introduction}

Many routes to the Riemann hypothesis translate zero location into positivity, self-adjointness,
or boundedness.  Li's criterion encodes the zeros through powers of the Cayley coordinate
\(1-\rho^{-1}\) \cite{Li1997,BombieriLagarias1999}; Keiper's generating series already centers
that coordinate at the arithmetic point \(s=1\) \cite{Keiper1992}; Voros showed how an off-line
zero creates exponential large-order behavior \cite{Voros2014}.  Hilbert--P\'olya and trace-formula
programs instead seek a positive spectral realization.  Recent work makes both the promise and
the missing positivity step especially explicit \cite{ConnesConsani2021,Connes2026,Hedenmalm2026}.

A separate body of evidence warns that fixed local observation is too weak.  Small
functional-equation-preserving perturbations can insert off-line zeros
\cite{GauthierXarles2009}.  In the constructions below, the warning is exact: a divisor quartet
moves to infinite height while every holomorphic jet converges locally uniformly, the
hypertranscendental differential stratum remains fixed, and a nonnegative charge density loses
all mass on bounded windows.  Yet its total charge does not change.

These facts are usually read as unrelated defects of particular methods.  The present paper gives
a common mechanism.  A marked source \(c\) determines the singular evolution
\[
  U_c(x;s):=\exp\!\left(\frac{x}{s-c}\right).
\]
If one evolves the meromorphic logarithmic response before extracting its residue, each divisor
point becomes one coherent exponential mode.  The derivative orders of that mode form a Poisson
packet.  At fixed \(x\), a high divisor has packet drift of order \(|\rho|^{-2}\); its nontrivial
orders therefore move beyond any fixed observation regime.  At the inverse-drift time
\(x=|\rho-c|^2\tau\), the radial factor cancels exactly.  This is not a limiting approximation.

The cancellation has a sharper complex form:
\[
  \frac{|\rho-c|^2\tau}{\rho-c}
  =\overline{\rho-c}\,\tau.
\]
The reciprocal Cayley coordinate is converted into direct conjugate displacement.  Consequently
spectral escape becomes rapid phase, not information loss.  After factorial summation of the
derivative orders, the raw energy \(e^{2\Rea(\rho-c)\tau}\) acquires the factor \(e^\tau\).
The neutral locus is therefore
\[
  2\Rea(\rho-c)+1=0.
\]
We call this line the \emph{source horizon}.

The Riemann case is distinguished twice.  First, the pole of the Euler factor places the
arithmetic source at \(c=1\).  Second,
\[
  1+\frac1{\rho-1}=\frac{\rho}{\rho-1}
  =\left(1-\frac1\rho\right)^{-1},
\]
so the source mode is exactly the Poisson completion of the inverse Li orbit.  The source horizon
is then \(\Rea\rho=\tfrac12\), and reflection \(\rho\mapsto1-\rho\) reverses the resulting
Lyapunov rate.

\subsection*{Main results}

The paper proves four closed statements.

\begin{enumerate}[leftmargin=2.2em]
\item The \emph{source inversion law}: canonical observation converts reciprocal displacement
into conjugate direct displacement.
\item The \emph{source-horizon conservation law}: the complete derivative-order energy is an
exact Poisson packet with rate \(2\Rea(\rho-c)+1\).
\item The \emph{Riemann source law}: the actual prime--Gamma logarithmic response transports every
\(\xi\)-zero residue into this packet, with no simplicity assumption, and the mode is the Poisson
completion of the inverse Li--Cayley orbit.
\item The \emph{escape--renormalization law}: explicit off-line quartets are simultaneously
invisible to every fixed compact holomorphic jet, to a complete rational differential matroid,
and to every bounded Ward window, but reappear exactly at canonical source time.
\end{enumerate}

The first two statements hold for every meromorphic function.  The last two identify why the
Riemann critical line is the neutral line selected by its arithmetic source and why several
apparently strong observers lose the same off-line data.

\subsection*{Novelty boundary}

No novelty is claimed for logarithmic residues, coherent-state Poisson statistics
\cite{Bargmann1961,Glauber1963}, Li's criterion, exponential Li asymptotics, or symmetric
polynomial perturbations.  The contribution is the exact conjunction: source evolution is
performed before residue extraction; its mode is identified with the factorial completion of the
Li orbit; inverse-drift time converts reciprocal to direct spectral displacement; the critical
line becomes the attenuation--order-entropy balance line; and the same law resolves the
compact-open, differential-algebraic, and Ward escape examples.  We found no prior source stating
this source-horizon theorem or its Riemann specialization.  That database statement is evidence,
not a claim that literature search proves absence.

\section{Meromorphic source evolution}

Let \(F\) be meromorphic near \(\rho\).  We use the signed divisor order convention
\[
  F(s)=(s-\rho)^m h(s),\qquad h(\rho)\ne0,
\]
where \(m>0\) for a zero and \(m<0\) for a pole.  Its logarithmic response is
\[
  \cR_F(s):=\frac{F'(s)}{F(s)}.
\]

\begin{proposition}[Residue transport]\label{prop:transport}
Let \(c,\rho\in\C\) with \(\rho\ne c\), and let \(m=\ord_\rho F\).  Then, for every
\(x\in\C\),
\[
 \Res_{s=\rho}
   \left(\cR_F(s)e^{x/(s-c)}\right)
 =m e^{x/(\rho-c)}.
\]
No simplicity assumption is required.
\end{proposition}

\begin{proof}
The local factorization gives
\[
  \cR_F(s)=\frac{m}{s-\rho}+\frac{h'(s)}{h(s)}.
\]
The second term is holomorphic at \(\rho\), while \(e^{x/(s-c)}\) is holomorphic there because
\(\rho\ne c\).  Multiplying the principal part by the value of the holomorphic factor proves the
formula.
\end{proof}

\begin{definition}
The transported residue mode is
\[
  M_{m,c,\rho}(x):=m e^{x/(\rho-c)}.
\]
Write
\[
  b_{c,\rho}:=(\rho-c)^{-1},\qquad
  d_{c,\rho}:=|b_{c,\rho}|^2,\qquad
  \kappa_{c,\rho}:=|1+b_{c,\rho}|^2-1.
\]
\end{definition}

Every evolution derivative is explicit:
\begin{equation}\label{eq:jets}
  \partial_x^r M_{m,c,\rho}(x)
  =m b_{c,\rho}^{\,r}e^{b_{c,\rho}x}.
\end{equation}

\section{Poisson completion and order-space conservation}

The elementary exponential series becomes structurally useful when its geometric parameter is the
Cayley node.

\begin{proposition}[Poisson completion]\label{prop:poisson-completion}
For \(a,x\in\C\),
\[
 e^{-x}\sum_{n=0}^{\infty}\frac{x^n}{n!}a^n=e^{(a-1)x}.
\]
Consequently,
\[
 M_{m,c,\rho}(x)
 =m e^{-x}\sum_{n=0}^{\infty}\frac{x^n}{n!}
   \left(1+\frac1{\rho-c}\right)^n.
\]
\end{proposition}

\begin{proof}
Absolute convergence permits direct summation:
\[
 e^{-x}\sum_{n\ge0}\frac{(ax)^n}{n!}=e^{-x}e^{ax}=e^{(a-1)x}.
\]
Use \(a=1+(\rho-c)^{-1}\).
\end{proof}

For \(\lambda\ge0\), write
\[
  p_r(\lambda):=e^{-\lambda}\frac{\lambda^r}{r!}
  =\Pr\{\Pois(\lambda)=r\}.
\]

\begin{theorem}[Order-space conservation]\label{thm:order-conservation}
For \(x\ge0\) and \(r\in\N_0\),
\[
 \frac{x^r}{r!}\left|\partial_x^rM_{m,c,\rho}(x)\right|^2
 =
 |m|^2e^{\kappa_{c,\rho}x}\,
 p_r(d_{c,\rho}x).
\]
In particular,
\[
 \sum_{r=0}^{\infty}
 \frac{x^r}{r!}\left|\partial_x^rM_{m,c,\rho}(x)\right|^2
 =|m|^2e^{\kappa_{c,\rho}x}.
\]
\end{theorem}

\begin{proof}
From \cref{eq:jets},
\[
 \frac{x^r}{r!}|\partial_x^rM(x)|^2
 =|m|^2e^{2\Rea(b_{c,\rho})x}
   \frac{(d_{c,\rho}x)^r}{r!}.
\]
Since
\[
 \kappa_{c,\rho}=2\Rea(b_{c,\rho})+d_{c,\rho},
\]
inserting \(e^{d_{c,\rho}x}e^{-d_{c,\rho}x}=1\) gives the first identity.  Summing
the Poisson weights gives the second.
\end{proof}

The theorem separates two invariants:

\begin{itemize}[leftmargin=2em]
\item \(d_{c,\rho}\) is packet velocity in derivative order;
\item \(\kappa_{c,\rho}\) is total exponential rate.
\end{itemize}

Neither invariant requires a choice of phase.  Together they recover the spectral displacement up
to conjugation.

\begin{proposition}[Ridge reconstruction]\label{prop:ridge}
Put \(z=\rho-c\ne0\), \(d=|z|^{-2}\), and
\(\kappa=|1+z^{-1}|^2-1\).  Then
\[
  \Rea z=\frac12\left(\frac{\kappa}{d}-1\right),
  \qquad
  (\Ima z)^2=\frac1d-\frac14\left(\frac{\kappa}{d}-1\right)^2.
\]
\end{proposition}

\begin{proof}
The identity
\[
  \kappa=\frac{2\Rea z+1}{|z|^2}=(2\Rea z+1)d
\]
gives the real part.  Then \(|z|^2=d^{-1}=(\Rea z)^2+(\Ima z)^2\).
\end{proof}

\section{The source-horizon law}

\begin{definition}
For \(\tau\ge0\), define canonical source time
\[
  x_{c,\rho}(\tau):=|\rho-c|^2\tau
\]
and transverse source rate
\[
  \nu_c(\rho):=2\Rea(\rho-c)+1.
\]
The \emph{source horizon} is the neutral locus \(\nu_c(\rho)=0\).
\end{definition}

\begin{theorem}[Source inversion]\label{thm:inversion}
For every \(c,\rho\in\C\), \(\rho\ne c\), and \(\tau\ge0\),
\[
  \frac{x_{c,\rho}(\tau)}{\rho-c}
  =\overline{\rho-c}\,\tau.
\]
Consequently,
\[
  M_{m,c,\rho}(x_{c,\rho}(\tau))
  =m e^{\overline{\rho-c}\tau}.
\]
\end{theorem}

\begin{proof}
For \(z=\rho-c\ne0\), \(z^{-1}=\bar z/|z|^2\).  Multiplication by
\(|z|^2\tau\) gives the first identity, and the second follows from the definition of \(M\).
\end{proof}

\begin{theorem}[Source-horizon conservation]\label{thm:horizon}
For every \(r\in\N_0\),
\begin{equation}\label{eq:horizon}
 \frac{x_{c,\rho}(\tau)^r}{r!}
 \left|\partial_x^rM_{m,c,\rho}(x)
   \big|_{x=x_{c,\rho}(\tau)}\right|^2
 =
 |m|^2e^{\nu_c(\rho)\tau}p_r(\tau).
\end{equation}
Equivalently,
\begin{equation}\label{eq:entropy}
 \frac{x_{c,\rho}(\tau)^r}{r!}
 \left|\partial_x^rM_{m,c,\rho}(x)
   \big|_{x=x_{c,\rho}(\tau)}\right|^2
 =
 |m|^2e^{2\Rea(\rho-c)\tau}\frac{\tau^r}{r!}.
\end{equation}
The complete channel mass is
\[
 \sum_{r=0}^{\infty}
 \frac{x_{c,\rho}(\tau)^r}{r!}
 \left|\partial_x^rM_{m,c,\rho}(x)
   \big|_{x=x_{c,\rho}(\tau)}\right|^2
 =|m|^2e^{\nu_c(\rho)\tau}.
\]
\end{theorem}

\begin{proof}
At \(x=x_{c,\rho}(\tau)\),
\[
 d_{c,\rho}x=|\rho-c|^{-2}|\rho-c|^2\tau=\tau
\]
and
\[
 \kappa_{c,\rho}x
 =\frac{2\Rea(\rho-c)+1}{|\rho-c|^2}
   |\rho-c|^2\tau
 =\nu_c(\rho)\tau.
\]
Insert both identities into \cref{thm:order-conservation}.  Expanding
\(p_r(\tau)=e^{-\tau}\tau^r/r!\) gives \cref{eq:entropy}; summing gives the final formula.
\end{proof}

\begin{corollary}[Horizon classification]\label{cor:horizon-classification}
\[
\begin{aligned}
 \nu_c(\rho)&>0 &&\Longleftrightarrow&
   \Rea\rho&>\Rea c-\tfrac12,\\
 \nu_c(\rho)&=0 &&\Longleftrightarrow&
   \Rea\rho&=\Rea c-\tfrac12,\\
 \nu_c(\rho)&<0 &&\Longleftrightarrow&
   \Rea\rho&<\Rea c-\tfrac12.
\end{aligned}
\]
\end{corollary}

\begin{remark}[Attenuation--entropy balance]
Equation \eqref{eq:entropy} isolates the mechanism.  Canonical inversion contributes raw mode
energy \(e^{2\Rea(\rho-c)\tau}\).  The factorial order tower contributes
\(\sum_r\tau^r/r!=e^\tau\).  The neutral line is exactly where these two factors balance.
\end{remark}

\begin{corollary}[Vertical parabolic renormalization]\label{cor:vertical}
Let
\[
  \rho_T=\frac12+a+iT,\qquad T\ne0,\qquad c=1,
\]
and
\[
  x_T(\tau)=\left((a-\tfrac12)^2+T^2\right)\tau.
\]
Then
\[
 \frac{x_T(\tau)^r}{r!}
 \left|\partial_x^r
   \left(m e^{x/(\rho_T-1)}\right)_{x=x_T(\tau)}\right|^2
 =
 |m|^2e^{2a\tau}p_r(\tau).
\]
All dependence on \(T\) has disappeared from the energy packet.
\end{corollary}

\section{The arithmetic source and the Riemann divisor}

Let
\[
  \xi(s)=\frac12s(s-1)\pi^{-s/2}\Gamma(s/2)\zeta(s)
\]
and
\[
  \Gamma_{\R}(s):=\pi^{-s/2}\Gamma(s/2).
\]
Its logarithmic response is \(\cR_\xi=\xi'/\xi\).  On \(\Rea s>1\),
\begin{equation}\label{eq:arithmetic-source}
 \cR_\xi(s)
 =\frac1s+\frac1{s-1}
   \frac{\Gamma_{\R}'(s)}{\Gamma_{\R}(s)}
   -\sum_{n=1}^{\infty}\frac{\Lambda(n)}{n^s}.
\end{equation}
Thus the marked point \(c=1\) is already present in the completed prime--Gamma source; it is not
chosen from the zero set.

Define
\[
  \cR_{\xi,x}(s):=\cR_\xi(s)e^{x/(s-1)}.
\]

\begin{theorem}[Riemann source transport]\label{thm:riemann-transport}
If \(\rho\) is a nontrivial zero of \(\xi\) of analytic multiplicity \(m_\rho\), then
\[
 \Res_{s=\rho}\cR_{\xi,x}(s)
 =m_\rho e^{x/(\rho-1)}.
\]
All source-horizon identities apply with \(c=1\) and \(m=m_\rho\).
\end{theorem}

\begin{proof}
Nontrivial zeros satisfy \(\rho\ne1\).  Apply \cref{prop:transport} to \(F=\xi\).
\end{proof}

\subsection{Exact relation to the Li orbit}

Put
\[
  w_\rho:=1-\frac1\rho=\frac{\rho-1}{\rho},
  \qquad
  a_\rho:=w_\rho^{-1}=\frac{\rho}{\rho-1}.
\]
Then \(a_\rho-1=(\rho-1)^{-1}\).  Hence
\begin{equation}\label{eq:li-poisson}
 e^{x/(\rho-1)}
 =e^{-x}\sum_{n=0}^{\infty}\frac{x^n}{n!}a_\rho^n.
\end{equation}
Li's coefficients use the discrete powers \(w_\rho^n\)
\cite{Li1997,BombieriLagarias1999}; false RH produces exponentially dominant Cayley powers
\cite{Voros2014}.  Equation \eqref{eq:li-poisson} identifies the continuous source residue with
the exact Poisson completion of the inverse orbit \(w_\rho^{-n}\).  Discrete Cayley dynamics,
continuous source evolution, and derivative-order transport are therefore three levels of one
semigroup.

For \(c=1\),
\[
  d_\rho=|\rho-1|^{-2},
  \qquad
  \kappa_\rho=(2\Rea\rho-1)d_\rho.
\]
The source horizon is
\[
  \Rea\rho=\frac12.
\]
Moreover,
\begin{equation}\label{eq:riemann-reconstruction}
 \Rea\rho
 =\frac12\left(1+\frac{\kappa_\rho}{d_\rho}\right),
 \qquad
 (\Ima\rho)^2
 =\frac1{d_\rho}
  -\frac14\left(1-\frac{\kappa_\rho}{d_\rho}\right)^2.
\end{equation}

\begin{corollary}[Riemann Lyapunov orientation]\label{cor:riemann-rate}
Every strict-right zero has positive source rate
\[
  \nu_1(\rho)=2\Rea\rho-1>0.
\]
The reflected zero \(1-\rho\) has rate \(-\nu_1(\rho)\), while conjugation reverses phase and
preserves rate.  Thus an off-line quartet is a stable--unstable, phase-reversed source packet,
whereas a critical-line pair is neutral.
\end{corollary}

\begin{proof}
The sign statement is \cref{cor:horizon-classification}.  Direct calculation gives
\[
 \nu_1(1-\rho)=2\Rea(-\rho)+1=1-2\Rea\rho=-\nu_1(\rho),
\]
and \(\Rea\bar\rho=\Rea\rho\).
\end{proof}

\section{Fixed-scale blindness and exact escape}

We now exhibit a family for which compact analysis, differential algebra, and bounded-window
charge all lose an off-line divisor, while \cref{thm:horizon} retains it.

Fix \(0<a<\tfrac12\), put \(T_n=n+1\), and define
\[
 G_{a,n}(s):=
 \left[1+T_n^{-2}\left((s-\tfrac12)-a\right)^2\right]
 \left[1+T_n^{-2}\left((s-\tfrac12)+a\right)^2\right],
\]
\[
 F_{a,n}(s):=G_{a,n}(s)\xi(s),
 \qquad
 \rho_{a,n}:=\frac12+a+iT_n.
\]
The normalization makes \(G_{a,n}\to1\) locally uniformly.  It also preserves
\(s\mapsto1-s\) and Schwarz conjugation.

\begin{theorem}[Compact-open hypertranscendental escape]\label{thm:escape}
For every \(n\):

\begin{enumerate}[label=\textup{(\roman*)},leftmargin=2.4em]
\item \(F_{a,n}\) is entire and hypertranscendental over \(\C(s)\);
\item \(F_{a,n}(1-s)=F_{a,n}(s)\) and
  \(F_{a,n}(\bar s)=\overline{F_{a,n}(s)}\);
\item \(F_{a,n}(\rho_{a,n})=0\) and
  \(\Rea\rho_{a,n}\ne\tfrac12\);
\item for every \(j\ge0\),
\[
  \partial_s^jF_{a,n}\longrightarrow\partial_s^j\xi
\]
locally uniformly on \(\C\).
\end{enumerate}

The same statements hold with \(G_{a,n}\) replaced by \(G_{a,n}^N\) for every positive
integer \(N\).
\end{theorem}

\begin{proof}
The symmetry identities follow by expanding the two paired quadratic factors.  At
\(\rho_{a,n}\), the first bracket equals \(1+T_n^{-2}(iT_n)^2=0\).  Local uniform convergence
of \(G_{a,n}\) and all its derivatives to \(1\) follows from \(T_n^{-2}\to0\); multiplication by
the fixed entire function \(\xi\), followed by the Weierstrass theorem for derivatives, gives
(iv).  Multiplication by a nonzero polynomial preserves hypertranscendence: a differential
algebraic relation for \(G_{a,n}\xi\) over \(\C(s)\) would divide by the rational polynomial and
give one for \(\xi\).  The power \(G_{a,n}^N\) is treated identically.
\end{proof}

\subsection{Differential-matroid blindness}

For a nonzero meromorphic \(H\), define its even-shift logarithmic packet
\[
  \Delta_m(H)(s)
  :=\frac{H'(s+2m)}{H(s+2m)}-\frac{H'(s)}{H(s)},
  \qquad m\in\N_0.
\]
Since
\[
 \Delta_m(F_{a,n})
 =\Delta_m(\xi)+\Delta_m(G_{a,n})
\]
and \(\Delta_m(G_{a,n})\in\C(s)\), every differential dependence matroid of the packet over
\(\C(s)\) is unchanged by the inserted quartet.  The same remains true after adjoining every
scalar Euler-valuation layer used by the prime source, and after replacing \(G_{a,n}\) by
\(G_{a,n}^N\).

\begin{theorem}[Exact Kolchin-stratum collision]\label{thm:matroid}
For every index set \(A\), grading \(g:A\to\N_0\), \(n\), and \(N\ge1\), the represented
differential matroid of
\[
  \bigl(\Delta_{g(q)}(G_{a,n}^N\xi)\bigr)_{q\in A}
\]
over the rational coordinate differential field is exactly the represented differential matroid
of
\[
  \bigl(\Delta_{g(q)}(\xi)\bigr)_{q\in A}.
\]
The equality persists over the differential field obtained by adjoining all scalar
prime-valuation layers.
\end{theorem}

\begin{proof}
For each \(q\), the two family elements differ by
\(\Delta_{g(q)}(G_{a,n}^N)\), an element of the base differential field.  Translating each
coordinate by a base-field element is an invertible affine change over that base.  It preserves
all differential-algebraic dependence subsets, hence the represented matroid.  The same argument
applies after any base-field enlargement containing the correction.
\end{proof}

\subsection{Ward charge escape}

The quartet has the nonnegative density
\begin{equation}\label{eq:ward-density}
 \cW_{a,T}(t)
 :=
 2\left[
 \frac{2(a+\tfrac12)}
 {(a+\tfrac12)^2+(t+T)^2}
 +
 \frac{2(a+\tfrac12)}
 {(a+\tfrac12)^2+(t-T)^2}
 \right].
\end{equation}

\begin{theorem}[Quantized Ward escape]\label{thm:ward}
For every \(T>0\),
\[
  \cW_{a,T}(t)\ge0,
  \qquad
  \int_{\R}\cW_{a,T}(t)\dd t=8\pi.
\]
For every \(R\ge0\),
\[
 \sup_{|t|\le R}\cW_{a,T_n}(t)\longrightarrow0.
\]
For quartet multiplicity \(N\), the density and total mass are multiplied by \(N\), so the
mass is \(8\pi N\) although every bounded window still loses it.
\end{theorem}

\begin{proof}
Nonnegativity is immediate from \eqref{eq:ward-density}.  Each translated Lorentzian satisfies
\[
 \int_{\R}
 \frac{2(a+\tfrac12)}
 {(a+\tfrac12)^2+(t\pm T)^2}\dd t=2\pi;
\]
the outer factor \(2\) therefore gives \(8\pi\).  If \(|t|\le R<T\), both denominators are at
least \((T-R)^2\), and
\[
 0\le\cW_{a,T}(t)\le\frac8{(T-R)^2}.
\]
This bound proves uniform convergence on the window.  Multiplicity is additive in the logarithmic
response.
\end{proof}

\subsection{The missing information is restored}

The zero \(\rho_{a,n}\) escapes vertically.  Let \(m_{a,n}\) be its actual analytic multiplicity
in \(F_{a,n}\); no simplicity assumption is needed.  At source \(c=1\),
\[
 x_{a,n}(\tau)
 =|\rho_{a,n}-1|^2\tau
 =\left((a-\tfrac12)^2+T_n^2\right)\tau.
\]

\begin{corollary}[Escape--renormalization]\label{cor:escape-renormalization}
For every \(r\in\N_0\) and \(\tau\ge0\),
\[
 \frac{x_{a,n}(\tau)^r}{r!}
 \left|
 \partial_x^r
 \left(m_{a,n}e^{x/(\rho_{a,n}-1)}\right)_
   {x=x_{a,n}(\tau)}
 \right|^2
 =
 |m_{a,n}|^2e^{2a\tau}p_r(\tau).
\]
The packet shape and exponential rate are exactly independent of height.
\end{corollary}

\begin{proof}
Apply \cref{cor:vertical} with \(T=T_n\) and \(m=m_{a,n}\).
\end{proof}

\Cref{thm:escape,thm:matroid,thm:ward,cor:escape-renormalization} give one precise
nonproperness law.  The following data can remain fixed or converge:

\begin{itemize}[leftmargin=2em]
\item every compact holomorphic jet;
\item functional equation and Schwarz symmetry;
\item hypertranscendence and the complete rational differential matroid;
\item every bounded Ward window.
\end{itemize}

Nevertheless the divisor is off line at every stage, its Ward charge is fixed, and its
source-horizon packet is recovered exactly after canonical rescaling.  Fixed observers do not
contradict the conserved observer; they are its non-co-moving projections.

\section{Consequences for Riemann-hypothesis routes}

\subsection{Why squaring and fixed scale lose orientation}

The energy packet \eqref{eq:horizon} retains the sign of
\(\nu_1(\rho)=2\Rea\rho-1\) in its total rate.  Phase-insensitive squaring does not itself erase
this orientation; observing at a fixed scale does.  As \(|\Ima\rho|\to\infty\),
\[
  d_\rho=|\rho-1|^{-2}\to0.
\]
At fixed \(x\), the Poisson mean \(d_\rho x\) collapses to zero and cannot resolve a moving
derivative-order ridge.  At \(x=\tau/d_\rho\), the ridge is restored at mean \(\tau\) and the
rate becomes \(2\Rea\rho-1\).  Thus the correct distinction is not ``squared versus unsquared''
alone, but ``fixed observer versus source-co-moving observer.''

\subsection{What the \(8\pi\) index does and does not say}

The Ward law counts a quartet and survives escape, but it supplies no upper bound forcing the
count to vanish.  Indeed, \cref{thm:ward} permits arbitrary mass \(8\pi N\) in the identical
compact-open and differential stratum.  The source-horizon law does not manufacture such a bound.
It answers a different closed question: it identifies the exact oriented dynamic variable that
the static index omits.  Ward mass records \emph{how much} divisor charge escapes; the Lyapunov
rate records \emph{to which side} of the source horizon it belongs.

\subsection{The new route boundary}

For the completed Riemann response, RH is equivalent to neutral source rate for every nontrivial
zero:
\[
  \mathrm{RH}
  \quad\Longleftrightarrow\quad
  \nu_1(\rho)=0\ \text{for every nontrivial zero }\rho.
\]
Taken alone, this equivalence is not a proof and receives no novelty credit.  The advance is the
source-derived realization behind it:

\begin{enumerate}[leftmargin=2.2em]
\item \(\nu_1(\rho)\) is the exact canonical-time growth rate of an actual residue of the
prime--Gamma response;
\item the same residue is the Poisson completion of the inverse Li orbit;
\item its packet drift and rate reconstruct the zero;
\item explicit countermodels prove why every fixed compact, finite differential, or bounded-window
observer can miss it;
\item inverse-drift renormalization recovers it without a limiting error.
\end{enumerate}

This replaces the search for another fixed positive scalar by a sharply specified dynamical
problem: control the source-generated semigroup at its canonical, divisor-dependent scale.  The
present paper closes the information-loss problem and identifies the oriented variable.  It does
not claim the additional arithmetic estimate that would exclude positive-rate modes.

\section{Relation to prior work}

Li proved positivity of his coefficient sequence equivalent to RH
\cite{Li1997}; Bombieri and Lagarias placed the criterion in a general multiset framework and
related it to Weil positivity \cite{BombieriLagarias1999}.  Generalized Li criteria move the
basepoint or target zero-free subregions \cite{Freitas2007,Sekatskii2014}.  The present law differs
in both operation and output: it evolves the logarithmic response before taking a local residue,
Poisson-completes the inverse orbit, and then chooses the exact inverse-drift time at which the
reciprocal spectral variable becomes direct conjugate displacement.

Voros's analysis gives the closest asymptotic neighbor: off-line zeros create exponential
large-order Li behavior \cite{Voros2014}.  Here the exponential rate is obtained channel by
channel through an exact identity, and the vertical height cancels at finite rescaled time rather
than only in a large-order asymptotic.

Coherent exponentials and their Poisson occupation statistics are classical
\cite{Bargmann1961,Glauber1963}.  Their role here is not a quantum analogy: the coherent mode is
the literal residue of the evolved arithmetic response.  The source inversion
\(z^{-1}|z|^2\tau=\bar z\tau\) and the resulting attenuation--order-entropy horizon are then exact
properties of meromorphic divisor transport.

Gauthier and Xarles showed broad instability of RH under small functional-equation-preserving
perturbations \cite{GauthierXarles2009}.  We assign no novelty to quartet insertion or
compact-open instability alone.  The new statement is their exact collision with every
holomorphic jet, a complete differential matroid, quantized Ward escape, and source-horizon
renormalization in one explicit family.

Connes and Consani analyzed scaling and the limitations of a positivity attempt in a semilocal
operator framework \cite{ConnesConsani2021}; Connes's recent survey stresses the continuing need
to connect finite arithmetic approximations to the full divisor \cite{Connes2026}.  Hedenmalm's
recent theta differential problem likewise isolates a missing self-adjoint structure
\cite{Hedenmalm2026}.  Our theorem is neither a no-go for those full programs nor a substitute
Hilbert space.  It specifies the noncompact, scale-covariant divisor datum that any successful
global carrier must retain.

\section{Closed scope}

The paper answers the question it poses:

\begin{quote}
When a divisor point becomes invisible to compact analytic observation, finite differential
algebra, and bounded charge windows, its information is transported into a source-generated
coherent mode.  Radial displacement fixes observation time, imaginary displacement becomes phase,
and signed distance from the source horizon becomes the exact Lyapunov rate.
\end{quote}

For the Riemann response, the source is \(1\), the horizon is the critical line, and the coherent
mode is the Poisson completion of the inverse Li orbit.  The conclusion is unconditional and
contains no simplicity, rightmost-zero, finite-off-line-divisor, or RH assumption.

The theorem does not assert that all Riemann modes are neutral.  It proves what the existing
blindness and charge laws are projections of, why the critical half-line is selected by the
arithmetic source, and which exact co-moving observable retains the orientation that fixed
scalarizations lose.

\begin{thebibliography}{99}

\bibitem{Bargmann1961}
V.~Bargmann,
\newblock On a Hilbert space of analytic functions and an associated integral transform. I,
\newblock \emph{Comm. Pure Appl. Math.} \textbf{14} (1961), 187--214.
\newblock \url{https://doi.org/10.1002/cpa.3160140303}.

\bibitem{BombieriLagarias1999}
E.~Bombieri and J.~C. Lagarias,
\newblock Complements to Li's criterion for the Riemann hypothesis,
\newblock \emph{J. Number Theory} \textbf{77} (1999), 274--287.
\newblock \url{https://doi.org/10.1006/jnth.1999.2392}.

\bibitem{Connes2026}
A.~Connes,
\newblock The Riemann hypothesis: past, present and a letter through time,
\newblock arXiv:2602.04022 (2026).
\newblock \url{https://arxiv.org/abs/2602.04022}.

\bibitem{ConnesConsani2021}
A.~Connes and C.~Consani,
\newblock The scaling Hamiltonian,
\newblock \emph{J. Operator Theory} \textbf{85} (2021), 257--276.
\newblock \url{https://arxiv.org/abs/1910.14368}.

\bibitem{Freitas2007}
P.~Freitas,
\newblock A Li-type criterion for zero-free half-planes of Riemann's zeta function,
\newblock \emph{J. London Math. Soc.} \textbf{75} (2007), 399--414.
\newblock \url{https://arxiv.org/abs/math/0507368}.

\bibitem{GauthierXarles2009}
P.~M. Gauthier and X.~Xarles,
\newblock Perturbations of \(L\)-functions with or without non-trivial zeros off the critical line,
\newblock arXiv:0911.5135 (2009).
\newblock \url{https://arxiv.org/abs/0911.5135}.

\bibitem{Glauber1963}
R.~J. Glauber,
\newblock Coherent and incoherent states of the radiation field,
\newblock \emph{Phys. Rev.} \textbf{131} (1963), 2766--2788.
\newblock \url{https://doi.org/10.1103/PhysRev.131.2766}.

\bibitem{Hedenmalm2026}
H.~Hedenmalm,
\newblock Spectral interpretation of Riemann zeta zeros,
\newblock arXiv:2606.17494 (2026).
\newblock \url{https://arxiv.org/abs/2606.17494}.

\bibitem{Keiper1992}
J.~B. Keiper,
\newblock Power series expansions of Riemann's \(\xi\) function,
\newblock \emph{Math. Comp.} \textbf{58} (1992), 765--773.
\newblock \url{https://doi.org/10.1090/S0025-5718-1992-1122072-5}.

\bibitem{Li1997}
X.-J. Li,
\newblock The positivity of a sequence of numbers and the Riemann hypothesis,
\newblock \emph{J. Number Theory} \textbf{65} (1997), 325--333.
\newblock \url{https://doi.org/10.1006/jnth.1997.2137}.

\bibitem{Sekatskii2014}
S.~K. Sekatskii,
\newblock Generalized Bombieri--Lagarias theorem and generalized Li criterion with its arithmetic
interpretation,
\newblock \emph{Ukrainian Math. J.} \textbf{66} (2014), 371--383.
\newblock \url{https://doi.org/10.1007/s11253-014-0935-8}.

\bibitem{Voros2014}
A.~Voros,
\newblock Zeta functions over zeros of zeta functions and an exponential-asymptotic view of the
Riemann hypothesis,
\newblock arXiv:1403.4558 (2014).
\newblock \url{https://arxiv.org/abs/1403.4558}.

\end{thebibliography}

\end{document}
